%%%%%%%%%%  scan idx 38  |  folio 25  %%%%%%%%%%
Variantes de l'exemple précédent~: $G$ est un groupe algébrique défini sur $k$ et l'on ne considère que les représentations rationnelles de $G$ définies sur $k$~; on peut considérer le cas d'un groupe topologique, ou de Lie, et des représentations continues, etc\dots\footnote{On peut construire d'autres exemples de $\lambda$-anneaux, par exemple l'ensemble des classes d'équivalence de formes quadratiques, l'ensemble des classes de représentations d'une algébre de Lie, etc\dots Nous promettons que nous n'aurons pas besoin de tous ces exemples pour traiter le théorème de Riemann--Roch~!}

2) Soit $X$ une variété algébrique, irréductible ou non, définie sur le corps $k$. On pourra considérer le groupe quotient du groupe libre engendré par les classes de fibrés vectoriels sur $X$ (définis sur $k$), par le sous-groupe engendré par les éléments de la forme $(E\oplus E')-(E)-(E')$, sur lequel les opérations de $\lambda$-anneau seront définies par le produit tensoriel et les puissances extérieures. Si $X$ est complète, c'est un groupe libre engendré par les fibrés vectoriels indécomposables sur $X$ (grâce à Krull--Remak--Schmidt)~; la connaissance de cet anneau et de sa base équivaut à la classification complète des fibrés vectoriels sur $X$ et la connaissance des opérations tensorielles élémentaires sur ces fibrés. Si $k$ est algébriquement clos de caractéristique $0$, toutes les autres opérations tensorielles sur des classes de fibrés sont alors connues, à l'aide de ce qui va suivre dans 3). On peut faire un passage au quotient plus draconien, par le groupe engendré par les éléments $(E)-(E')-(E'')$ où $E$ est un fibré vectoriel extension de $E'$ par $E''$~; on trouve alors un $\lambda$-anneau, noté $\underline{K}(\xi(X))$, où $\xi(X)$ désigne la catégorie additive des fibrés vectoriels sur $X$, définis sur $k$. En utilisant les résultats de l'exemple 1), on montre que si $k$ est algébriquement clos\footnote{Condition inutile, cf.\ note de bas de page précédente.}, le $\lambda$-anneau $\underline{K}(\xi(X))$ est spécial, et qu'il en est de même du
